%! Solving Radical Equations \& Extraneous Solutions — Unit 6, Lesson 6.5 — Warm-Up (Answer Key)
\documentclass[10pt]{article}
\usepackage{algebra2-article}
\usepackage{algebra2-key}   % pulls in -boxes; do NOT also load -boxes
\newcommand{\TallMath}[1]{$\displaystyle #1\rule[-1.4em]{0pt}{3.2em}$}
% In-formula answer slot (\ans is text-mode and must never appear inside $...$).
\newcommand{\mans}[1]{{\color{keyred}\mathbf{#1}}}

\begin{document}

\pageheader{Unit 6, Lesson 6.5}{Warm-Up --- Answer Key}
\namedateperiod

\vspace{0.10in}

\begin{notesbox}{Getting Ready: Skills We'll Use Today}
\small Three short items. The third one is a trap, and you will meet it again in about ten minutes.

\vspace{4pt}
\textbf{1. Undoing a radical (Lessons 6.1 and 6.3).} Simplify each expression. Assume $x$ is large
enough that every radical is real.
\begin{center}
\renewcommand{\arraystretch}{1.7}
\begin{tabular}{@{}l l l@{}}
(a) \; $\left(\sqrt{x-3}\,\right)^{2}=$ \ans{$x-3$} &
(b) \; $\left(\sqrt[3]{2x}\,\right)^{3}=$ \ans{$2x$} &
(c) \; $\left(\sqrt{x}+3\right)^{2}=$ \ans{$x+6\sqrt{x}+9$} \\
\end{tabular}
\end{center}
\vspace{-4pt}
Parts (a) and (c) both squared something containing $\sqrt{x}$. One of them came out clean and the
other came out \emph{worse}. What was different about the one that came out clean?
\ansline{In (a) the radical was \emph{alone} --- squaring it erased it. In (c) the radical was part
of a sum, so squaring meant FOIL (Lesson 6.3), and the middle term $6\sqrt{x}$ still has a radical
in it.}

\vspace{0.30in}

\textbf{2. A one-way street.} Answer each question \textbf{yes} or \textbf{no}.
\begin{center}
\renewcommand{\arraystretch}{1.6}
\begin{tabularx}{0.96\linewidth}{@{}X c@{}}
\hline
If you know that $x=-4$, does it follow that $x^{2}=16$? & \ans{yes} \\
If you know that $x^{2}=16$, does it follow that $x=4$? & \ans{no ($x$ could be $-4$)} \\
\hline
\end{tabularx}
\end{center}
\vspace{-4pt}
Finish the sentence. \emph{Squaring both sides of a true equation always gives another true
equation, but going backwards} \ans{is not guaranteed --- the squared equation can be true for
values the original one was never true for.}

\vspace{4pt}
\textbf{3. Factoring, from Unit 3.} Solve by factoring.
\[
  x^{2}-11x+18=0 \qquad\Longrightarrow\qquad x=\mans{2} \ \text{ or } \ x=\mans{9}
\]
\vspace{-10pt}

\vspace{0.55in}

Hold on to those two numbers. Today an equation will hand you \emph{both} of them, and only one of
them will be telling the truth.

\end{notesbox}

\begin{teachernote}
\small Five minutes, no more. Every item is deliberately pre-assembled for a specific moment in the
notes.
\begin{itemize}[leftmargin=1.3em, nosep]
  \item \textbf{Item 1} is the reason Step 1 of today's method is \emph{isolate the radical}. Do not
        explain that yet --- just make sure both answers are on the board, side by side. When a
        student later squares $\sqrt{x}+2=5$ into $x+4=25$, point back here.
  \item \textbf{Item 2 is the entire lesson in two questions.} Squaring is a valid move but a
        \emph{one-way} one: it preserves every solution and can invent new ones. Get the phrase
        ``it can add answers'' said out loud now; the whole idea of an extraneous solution is
        already in the room.
  \item \textbf{Item 3} is the quadratic that $\sqrt{x+7}=x-5$ produces in Section 2. Students will
        recognize $2$ and $9$ when they appear, which makes the discovery that $2$ fails land much
        harder than if the numbers were new.
\end{itemize}
\end{teachernote}

\end{document}
